\documentclass{article}
\usepackage[UTF8]{ctex}
\usepackage{amsmath,amssymb,amsfonts}
\usepackage{enumitem}
\usepackage{geometry}
\geometry{a4paper, margin=1in}
% === 错题本排版优化（不参与题目标识） ===
\setlength{\parindent}{0pt}
\setlength{\parskip}{0.35em plus 0.10em minus 0.08em}
\setlength{\emergencystretch}{3em}
\setlist{itemsep=0.25em, topsep=0.40em, parsep=0pt}
\allowdisplaybreaks[2]
% === 错题本排版优化结束 ===
\begin{document}
\section*{北京师范大学2024年数学分析真题}
\section*{1．北京师范大学2024年数学分析试卷}

\subsection{题目1}
1．设 \(a_{1}>0, a_{n+1}=\ln \left(1+a_{n}\right)(n=1,2, \cdots)\) ，求 \(\lim _{n \rightarrow \infty} a_{n}\) 及 \(\lim _{n \rightarrow \infty} n a_{n}\) ．

\subsection{题目2}
2．设 \(f(x)\) 在 \([a,+\infty)\) 上连续，且当 \(x \rightarrow+\infty\) 时有渐近线 \(y=b x+c\) ，证明：\(f(x)\) 在 \([a,+\infty)\) 上一致连续。

\subsection{题目3}
3．证明：方程 \(3^{x}+7^{x}=2 \cdot 5^{x}\) 的实根只有 0 和 1 。

\subsection{题目4}
4．设 \(f(x)\) 在 \([0,1]\) 上连续且恒正，证明：对任意的正整数 \(n\) ，存在 \(\xi_{n} \in(0,1)\) ，使得

\[
\frac{1}{n} \int_{0}^{1} f(x) \mathrm{d} x=\int_{0}^{\xi_{n}} f(x) \mathrm{d} x+\int_{1-\xi_{n}}^{1} f(x) \mathrm{d} x
\]

且 \(\lim _{n \rightarrow \infty} n \xi_{n}=\frac{\int_{0}^{1} f(x) \mathrm{d} x}{f(0)+f(1)}\) ．

\subsection{题目5}
5．设 \(f(x)\) 是以 \(2 \pi\) 为周期的连续函数，记 \(f_{n}(x)=n \int_{x}^{x+\frac{1}{n}} f(t) \mathrm{d} t\) ，证明：\(\left\{f_{n}(x)\right\}\) 为连续可微的 \(2 \pi\)为周期的函数列，且在 \(\mathbb{R}\) 上一致收敛于 \(f(x)\) 。

\subsection{题目6}
6．判断级数 \(\sum_{n=2}^{\infty} \ln \left(1+\frac{(-1)^{n}}{n^{p}}\right)\) 的敛散性，\(p>0\) 。

\subsection{题目7}
7．计算极限 \(\lim _{n \rightarrow \infty} \int_{0}^{1} \frac{\sin ^{2} n x}{1+x^{2}} \mathrm{~d} x\) 。

\subsection{题目8}
8．设 \(f(x, y, z)\) 有连续的偏导数，且满足 \(-y \frac{\partial f}{\partial x}+x \frac{\partial f}{\partial y}+b \frac{\partial f}{\partial z} \geq c\) ，其中 \(b, c\) 为正实数，证明：当 \((x, y, z)\)沿着 \(x=b \cos t, y=b \sin t, z=b t, t \in[0,+\infty)\) 趋于无穷时，\(f(x, y, z) \rightarrow+\infty\) ．

\subsection{题目9}
9．设 \(\Omega\) 是 \(z=x^{2}+y^{2}\) 介于 \(z=0\) 和 \(z=2\) 之间的曲面，求曲面积分

\[
\iint_{\Sigma}\left(x^{2} \cos \alpha+y^{2} \cos \beta+z^{2} \cos \gamma\right) \mathrm{d} S .
\]

其中 \((\cos \alpha, \cos \beta, \cos \gamma)\) 为 \(\sum\) 的单位外法向量．

\subsection{题目10}
10．计算曲线积分

\[
\oint_{L}\left(y^{2}-z^{2}\right) \mathrm{d} x+\left(z^{2}-x^{2}\right) \mathrm{d} y+\left(x^{2}-y^{2}\right) \mathrm{d} z .
\]

其中 \(L\) 为 \(x^{2}+y^{2}+z^{2}=4\) 与 \(x^{2}+y^{2}=4\) 的交线，从 \(z\) 轴正向看为逆时针方向。
\end{document}
